
This study set out to answer: when AI assistants match user communication style, do users engage more? The answer is nuanced. AI systems do exhibit formality accommodation (r=0.152), and accommodation is associated with engagement---but the relationship is non-linear.

Moderate accommodation is associated with highest continuation rates (71.4\%), while both high (65.9\%) and low (60.9\%) accommodation are associated with lower rates. This inverted-U pattern suggests that optimal accommodation may lie in a middle range.

The implications are tentative. For chatbot designers, the findings suggest calibrating accommodation rather than maximizing it may be worth considering. For theory, the findings indicate that CAT predictions may require qualification in human-AI contexts. However, these are correlational findings from a single dataset, and causal mechanisms remain unverified.

Future work should test these patterns across multiple datasets and model families, manipulate accommodation levels experimentally to establish causal relationships, and investigate the mechanism underlying the inverted-U pattern.

